\chapter{1982:Aaron Klug}

\label{chap:chemistry1982}

The Nobel Prize in Chemistry for the year 1982 was awarded to Aaron Klug.

\textbf{Motivation:}

for his development of crystallographic electron microscopy and his structural elucidation of biologically important nucleic acid-protein complexes

\section{Aaron Klug}

\subsection*{Early Life and Education}

Aaron Klug was born on 1926-08-11 in Zelva, Lithuania.

\subsection*{Career and Major Research}

Klug's family moved to South Africa when he was a child. He studied at the University of the Witwatersrand and received his doctorate from the University of Cambridge. He worked at Birkbeck College in London and then at the Laboratory of Molecular Biology in Cambridge, where he developed crystallographic electron microscopy and used it to study the structures of viruses and other biological assemblies.

\subsection*{Nobel Prize-winning Work}

The work for which Aaron Klug was awarded the Nobel Prize in Chemistry in 1982 was described by the Nobel Committee as: "for his development of crystallographic electron microscopy and his structural elucidation of biologically important nucleic acid-protein complexes".

Klug combined electron microscopy with mathematical analysis to reconstruct three-dimensional structures from images, and he used these methods to determine the structure of the tobacco mosaic virus and, with his colleagues, the nucleosome.

\subsection*{Significance and Impact}

His imaging methods allowed the structures of large biological assemblies to be determined at high resolution, opening a new field of structural biology.

\subsection*{Later Career and Honors}

Klug worked at the Laboratory of Molecular Biology in Cambridge for most of his career and was awarded the Order of Merit in 1995. He died on 2018-11-20 in Cambridge, England.

\subsection*{Legacy}

The techniques Klug pioneered are the forerunners of modern cryo-electron microscopy.

\subsection*{Modern Developments}

Klug's development of crystallographic electron microscopy made it possible to visualize the structures of viruses and large macromolecular assemblies. His methods evolved into modern cryo-electron microscopy, which now routinely determines the structures of proteins and viruses at atomic resolution and is central to the development of vaccines and drugs, including for COVID-19.

\subsection*{Selected Publications}

"From Macromolecules to Biological Assemblies" (Nobel Lecture, 1982)

"Structure of the Nucleosome Core Particle at 7 A Resolution" (Nature, 1984)

\clearpage

\section*{Chapter Summary}

The 1982 prize honored Aaron Klug for developing crystallographic electron microscopy and for determining the structures of nucleic acid-protein complexes. His methods founded modern structural biology of large assemblies.

